(1) 棱镜内的两折射角满足 \(r_1+r_2=A\)。两界面的偏向角分别为 \(i-r_1\) 与 \(e-r_2\)，因此总偏向 \(\delta=i+e-A\)。Snell 定律 \( \sin i=n\sin r_1\)、\(\sin e=n\sin r_2\) 限定各面能否存在实传播角。
(2) 在对称且可行的最小偏向光线中入射、出射光路互易，\(i=e\)，故内部角 \(r_1=r_2=A/2\)。Snell 定律给 \(\sin[(A+\delta_{\min})/2]=n\sin(A/2)\)，从而 \(\delta_{\min}=2\arcsin[n\sin(A/2)]-A\)；反过来 \(n=\sin[(A+\delta_{\min})/2]/\sin(A/2)\)，是常用的折射率测量式。
(3) 固定 \(A\) 对 \(n\) 微分，得 \(d\delta_{\min}/dn=2\sin(A/2)/\sqrt{1-n^2\sin^2(A/2)}>0\)。如果 \(n=n(\lambda)\)，则波长微差对应 \(\Delta\delta_{\min}\simeq[2\sin(A/2)/\sqrt{1-n^2\sin^2(A/2)}](dn/d\lambda)\Delta\lambda\)，其符号由材料正常或反常色散确定。
(4) 对称传播要求 \(n\sin(A/2)\le1\)，在等号处折射光处于掠射边界，微分公式分母趋零、不能再线性外推；实际最小偏向测量通常取严格小于号。若 \(n\to1\)，\(\delta_{\min}\to0\)；若 \(A\ll1\) 且 \(nA/2\ll1\)，\(\arcsin(nA/2)\simeq nA/2\)，故 \(\delta_{\min}\simeq(n-1)A\)。
